\subsection{Программная реализация}

Языком разработки был выбран С++ в силу его низкоуровневости и скорости, а также с намереньем в дальнейшим внедрить эти наработки в конечную реализацию графовой БД на С++.

Программа построена по модульному принципу и состоит из набора компонентов, отвечающих за хранение графовых данных, моделирование распределённой среды и выполнение алгоритмов оптимизации. Архитектура проекта предусматривает использование абстрактных интерфейсов для основных подсистем, что позволяет заменять конкретные реализации без изменения остальной части программного кода. Такой подход обеспечивает возможность проведения сравнительных экспериментов с различными алгоритмами и конфигурациями системы.

\subsubsection{Классы представления графа}

Основой системы являются классы представления графа. Они определяют структуры данных, используемые для хранения вершин, рёбер и связанной служебной информации. Поскольку проект ориентирован на работу с графами различной природы и масштаба, представление графовых объектов реализовано максимально универсальным образом.

\textbf{Класс ключа вершины}

Первым базовым элементом является класс ключа вершины, фрагмент реализации которого представлен на листинге \ref{impl-lst:nodekey}.

\begin{lstlisting}[language=C++, caption=Класс NodeKey, label=impl-lst:nodekey]
template <typename KeyType> 
class NodeKey {
public:
KeyType key_value;

NodeKey<KeyType>& operator=(const NodeKey<KeyType>& other);
friend bool operator<(const NodeKey<KeyType>& lhs, const NodeKey<KeyType>& rhs);
friend bool operator>(const NodeKey<KeyType>& lhs, const NodeKey<KeyType>& rhs);
friend bool operator==(const NodeKey<KeyType>& lhs, const NodeKey<KeyType>& rhs);
friend bool operator!=(const NodeKey<KeyType>& lhs, const NodeKey<KeyType>& rhs);
\end{lstlisting}

Класс инкапсулирует идентификатор вершины и используется в качестве единого механизма адресации графовых объектов во всех подсистемах приложения. Использование отдельного класса вместо непосредственного применения примитивных типов позволяет унифицировать работу с ключами и при необходимости расширять их функциональность без изменения остальной части системы.

Класс реализован как шаблонный, благодаря чему тип идентификатора не фиксируется на этапе разработки. В качестве ключа могут использоваться как целочисленные типы данных, так и строки, UUID или другие пользовательские типы, удовлетворяющие требованиям сравнения и копирования. Такой подход повышает универсальность программной реализации и позволяет применять её для различных предметных областей.

Поскольку объекты данного класса используются в стандартных контейнерах языка C++, в частности в ассоциативных структурах данных, необходимо определить операции сравнения между экземплярами класса. Для этого перегружены операторы отношения \texttt{<} и \texttt{>}, а также операторы проверки равенства и неравенства. Наличие данных операторов позволяет использовать объекты класса в качестве ключей контейнеров \texttt{std::map}, \texttt{std::set} и других структур, основанных на упорядоченном хранении элементов.

Класс является шаблонным, позволяя использовать любой тип данных как ключ. Также из-за особенностей языка и стандартных структур требуется перегрузить часть операторов.

\textbf{Класс аттрибута}

Следующим элементом модели графа является класс параметра, общий вид которого представлен на листинге \ref{impl-lst:parameter}.

Данный класс предназначен для унифицированного хранения произвольных атрибутов вершин и рёбер графа. Использование единого контейнера параметров позволяет не привязывать реализацию графа к конкретному набору свойств и при необходимости расширять состав хранимых данных без изменения структуры основных классов.

\begin{lstlisting}[language=C++, caption=Класс Parameter, label=impl-lst:parameter]
class Parameter{
private:
    std::vector<std::byte> data;
public:
    void set_data(const std::vector<std::byte>& _data);
    std::vector<std::byte> get_data();
};
\end{lstlisting}

Внутри класса данные хранятся в виде массива байтов, что обеспечивает возможность сохранения объектов различных типов в едином формате. Такой подход может и, забегая вперёд, будет использоваться для хранения служебной информации, необходимой алгоритмам обработки графа, в том числе данным, применяемым при выполнении поиска путей и других операций анализа графовой структуры.

По своей реализации класс является простым контейнером данных и предоставляет только методы записи и чтения хранимого значения. Это обеспечивает минимальную сложность реализации и независимость механизма хранения параметров от логики работы остальных компонентов системы.

\textbf{Класс вершины}

Основным классом представления графа является класс вершины, общий вид которого представлен на листинге \ref{impl-lst:node}.

\begin{lstlisting}[language=C++, caption=Класс Node, label=impl-lst:node]
template <typename KeyType>
class Node {
private:
    NodeKey<KeyType> key;
public:
    std::map<std::string, Parameter> parameters;
    std::map<NodeKey<KeyType>, Edge<KeyType>> edges;

    NodeKey<KeyType> get_key() const;

    void add_edge(Edge<KeyType> new_edge);
    void add_edges(std::map<NodeKey<KeyType>, Edge<KeyType>> new_edges);

    bool remove_edge_to(NodeKey<KeyType> neighbour_key);
    bool remove_edge(Edge<KeyType> removing_edge);

    Node& operator=(const Node& other);

    bool operator<(const Node<KeyType>& other) const;
};
\end{lstlisting}

Экземпляр данного класса содержит всю информацию, необходимую для хранения вершины и её непосредственного окружения. Каждая вершина обладает уникальным ключом, представленным ранее описанным классом идентификатора, который используется для её однозначной идентификации внутри графа.

Атрибуты вершины хранятся в ассоциативном контейнере, где каждому строковому имени параметра соответствует объект класса Parameter. Такой подход позволяет динамически задавать набор свойств вершины без необходимости изменять структуру класса при добавлении новых атрибутов.

Информация о связях с соседними вершинами хранится в виде индексированной коллекции исходящих рёбер. В качестве ключа используется идентификатор соседней вершины, что обеспечивает быстрый доступ к необходимому ребру и позволяет эффективно выполнять операции поиска, добавления и удаления связей.

Хранение исходящих рёбер непосредственно внутри вершины упрощает реализацию алгоритмов обхода графа и поиска путей, поскольку информация о соседях доступна без обращения к дополнительным структурам данных. Кроме того, такой подход позволяет передавать вершину вместе с её локальным окружением между шардами в рамках процедур перераспределения данных. При этом при желании хранить рёбра в другом виде без особых проблем можно оставить коллекцию пустой без лишних затрат.

Класс предоставляет набор методов для добавления и удаления рёбер, а также для получения идентификатора вершины. Дополнительно реализованы оператор присваивания и оператор сравнения, необходимые для корректной работы со стандартными контейнерами языка C++.

\textbf{Класс ребра}

Последним базовым элементом модели графа является класс ребра, общий вид которого представлен на листинге \ref{impl-lst:edge}.

\begin{lstlisting}[language=C++, caption=Класс Edge, label=impl-lst:edge]
template <typename KeyType>
class Edge {
private:
    float weight;
    bool directional;
    std::map<std::string, Parameter> parameters;
    NodeKey<KeyType> from;
    NodeKey<KeyType> to;
public:
    float get_weight() const;
    bool is_directional() const;
    const std::map<std::string, Parameter>& get_parameters() const;
    const NodeKey<KeyType>& get_from();
    const NodeKey<KeyType>& get_to() const;

    Edge& get_reverted();
    
    Edge& operator=(const Edge& other);
    bool operator<(const Edge& other) const;

    // doesn't check extra parameters
    bool operator==(const Edge& other);
    
    NodeKey<KeyType> get_other(const NodeKey<KeyType>& node) const;
    NodeKey<KeyType> get_other(NodeKey<KeyType>* node) const;
};
\end{lstlisting}

Класс предназначен для хранения информации о связи между двумя вершинами графа. Помимо идентификаторов начальной и конечной вершин, объект ребра содержит вес связи и набор дополнительных параметров, представленных в виде коллекции объектов класса Parameter. Это позволяет использовать рёбра как для описания структуры графа, так и для хранения информации, необходимой алгоритмам обработки данных.

Для хранения информации о соединяемых вершинах используются поля from и to, содержащие идентификаторы начала и конца ребра соответственно. В случае направленного графа данные поля определяют направление перехода между вершинами и используются алгоритмами обхода при формировании допустимых маршрутов. Для ненаправленных рёбер такое разделение носит преимущественно технический характер и необходимо для единообразного хранения данных. Использование идентификаторов вершин вместо прямых ссылок на объекты позволяет упростить сериализацию графа, передачу данных между шардами и операции перераспределения вершин в распределённом хранилище.

Особенностью реализации является поддержка как направленных, так и ненаправленных рёбер в рамках единого класса. Для этого каждое ребро содержит признак направленности, определяющий правила его интерпретации при обходе графа. Такой подход позволяет использовать одну и ту же структуру данных для различных типов графов без необходимости создания отдельных классов.

Для упрощения работы с графом класс предоставляет вспомогательные методы получения противоположной вершины ребра по известной вершине, а также метод формирования обратного представления ребра, либо же просто объекта с поменянными местами исходящей и входящей вершинами для ненаправленных рёбер, для которых понятия исходящей и входящей вершины не имеют смысла. Данные операции активно используются при обходе графа и поиске соседних вершин, позволяя избежать дополнительной логики обработки направленности связей в алгоритмах более высокого уровня.

Кроме того, класс содержит методы доступа к основным характеристикам ребра и набор операторов сравнения, необходимых для корректной работы со стандартными контейнерами и алгоритмами языка C++.

\subsubsection{Классы имитации распределённой системы}

\textbf{Класс имитации шины}

Центральным элементом подсистемы имитации распределённой среды является интерфейс шины, общий вид которого представлен на листинге \ref{impl-lst:ibus}. Интерфейс в данном случае не является средством упрощения модификации, а средством избегания циклических зависимостей.

Данный интерфейс определяет набор операций, необходимых для взаимодействия между экземплярами хранилищ. В отличие от других интерфейсов системы, его основное назначение заключается не в обеспечении возможности замены реализаций, а в устранении циклических зависимостей между компонентами. Хранилища используют шину для обмена информацией друг с другом, при этом сама шина также должна иметь доступ к хранилищам для маршрутизации запросов. Выделение абстрактного интерфейса позволяет разорвать данную зависимость и упростить структуру проекта.

\begin{lstlisting}[language=C++, caption=Интерфейс IBus, label=impl-lst:ibus]
template <typename KeyType> 
class IBus {
public:
    virtual Node<KeyType> request_node(const NodeKey<KeyType>& node) = 0;
    virtual int send_add_node(const Node<KeyType>& node) = 0;
    virtual bool send_add_node(const Node<KeyType>& node, int storage_id) = 0;
    virtual bool send_remove_node(const NodeKey<KeyType>& node) = 0;
    virtual bool send_remove_node(const NodeKey<KeyType>& node, int storage_id) = 0;
    virtual int ask_who_has(int asker_id, NodeKey<KeyType> key) = 0;
    virtual void announce_add(NodeKey<KeyType> key, int storage_id, std::set<Edge<KeyType>> edges) = 0;
    virtual void announce_remove(NodeKey<KeyType> key, int storage_id) = 0;
    virtual float get_streaming_euristics_change(const Node<KeyType>& node, int storage_id) = 0;
    // запрашивает у source вершины, соседствующие с target
    virtual std::set<Node<KeyType>> ask_neigbours_to_storage(int source, int target) = 0;
    // запрашивает у source рёбра, идущие в target
    virtual std::set<Edge<KeyType>> ask_edges_to_storage(int source, int target) = 0;
};
\end{lstlisting}

Интерфейс предоставляет методы для получения вершин из удалённых хранилищ, добавления и удаления вершин, а также определения местоположения конкретной вершины в распределённой системе. Для поддержания актуальной информации о размещении данных предусмотрены механизмы уведомления о появлении и удалении вершин на шардах.

Отдельная группа методов предназначена для поддержки алгоритмов потокового распределения и оптимизации. Они позволяют оценивать целесообразность размещения вершины на конкретном шарде, а также получать информацию о связях между шардами без непосредственного доступа к их внутренним структурам данных.

Кроме того, интерфейс содержит методы для получения сведений о межшардовых связях. С их помощью можно запрашивать вершины и рёбра, соединяющие различные хранилища, что используется при выполнении распределённых алгоритмов обхода графа и при анализе качества текущего распределения данных.


\textbf{Класс хранилища}

Следующим элементом подсистемы имитации распределённой среды является хранилище, общий вид интерфейса которого представлен на листинге \ref{impl-lst:istorage}. 

Как и в случае интерфейса шины, основным назначением данного интерфейса является устранение циклических зависимостей между компонентами системы. Хранилища должны взаимодействовать друг с другом через шину, а шина, в свою очередь, должна иметь возможность обращаться к отдельным хранилищам. Использование абстрактного интерфейса позволяет организовать такое взаимодействие без прямой зависимости реализаций друг от друга.

Интерфейс определяет полный набор операций, необходимых для работы распределённого хранилища. К ним относятся подключение к шине, добавление и удаление вершин, получение доступа к локальным данным, обработка уведомлений об изменениях в других шардах, а также предоставление статистической информации о состоянии хранилища.

Отдельные методы предназначены для поддержки алгоритмов потокового распределения. Они позволяют оценивать влияние размещения новой вершины на качество текущего распределения и получать сведения о связях между различными шардами. Кроме того, интерфейс предоставляет операции для получения информации о межшардовых рёбрах и вершинах, имеющих соседей в других хранилищах.

\begin{lstlisting}[language=C++, caption=Интерфейс IStorage, label=impl-lst:istorage]
template <typename KeyType>
class IStorage {
public:
virtual int get_id() const = 0;
virtual void connect_to_bus(IBus<KeyType>* _bus) = 0;

virtual void get_add_announcement(NodeKey<KeyType> key, int announcer_id, std::set<Edge<KeyType>> edges) = 0 ;
virtual void get_remove_announcement(NodeKey<KeyType> key, int announcer_id) = 0;

virtual std::optional<std::set<Edge<KeyType>>> add_node(const Node<KeyType>& node) = 0;
virtual std::optional<std::set<Edge<KeyType>>> add_node_and_announce(const Node<KeyType>& node) = 0;

virtual bool remove_node(const NodeKey<KeyType>& key) = 0;
virtual bool remove_node_and_announce(const NodeKey<KeyType>& key) = 0;
   
virtual Node<KeyType>* get_node(const NodeKey<KeyType>& key) = 0;
virtual bool has_node(const NodeKey<KeyType>& key) const = 0;
virtual const std::map<NodeKey<KeyType>, Node<KeyType>>& get_all_nodes() const = 0;

virtual size_t size() const = 0;
virtual size_t internal_edges_size() const = 0;

virtual float get_streaming_euristics_change(const Node<KeyType>& node, long total_edges)  = 0;

// Получить набор узлов, имеющих соседей в указанном хранилище (копии)
virtual std::set<Node<KeyType>> get_nodes_with_neighbors_in_storage(int target_storage_id) const = 0;
virtual std::set<Edge<KeyType>> get_all_edges_to_storage(int target_storage_id) const = 0;
};
\end{lstlisting}

Поскольку интерфейс не совсем удобен как основа для реализации различных алгоритмов размещения данных, в системе дополнительно присутствует базовый класс хранилища, фрагмент которого представлен на листинге \ref{impl-lst:storage}. Данный класс содержит реализацию большинства стандартных операций и выступает в качестве основы для специализированных модификаций хранилища. Таким образом любые модификации могут вносить изменения только в нужные для них методы.

Базовый класс хранит уникальный идентификатор шарда, коллекцию принадлежащих ему вершин и информацию о межшардовых связях. Для хранения локальных данных используется ассоциативный контейнер, отображающий идентификаторы вершин в соответствующие объекты графа. Межшардовые связи сохраняются отдельно и группируются по идентификаторам удалённых хранилищ, что позволяет быстро определять наличие связей между различными частями распределённого графа.

Также класс содержит указатель на шину обмена сообщениями и статистическую информацию о количестве внутренних рёбер. Наличие общей реализации позволяет сосредоточиться на модификации только специфических методов при разработке новых алгоритмов распределения или хранения данных, не дублируя базовую функциональность работы с графом и взаимодействия с другими компонентами системы.

\begin{lstlisting}[language=C++, caption=Базовый класс Storage, label=impl-lst:storage]
template <typename KeyType>
class Storage : public IStorage<KeyType> {
protected:
typedef Node<KeyType> StorageNode;
typedef NodeKey<KeyType> Key;
int storage_id;
std::map<Key, StorageNode> nodes;
std::map<int, std::map<Key, std::map<Key, Edge<KeyType>>>> external_edges;
size_t internal_edges = 0;
IBus<KeyType> *bus = nullptr;
/* Простейшие реализациии методов интерфейса */
};

\end{lstlisting}

\textbf{Класс мастера}

Последним элементом подсистемы имитации распределённой среды является мастер, общий вид класса которого представлен на листинге \ref{impl-lst:master}. 

Мастер не хранит графовые данные напрямую, а оперирует абстракциями хранилищ, объединяя их в единую распределённую систему. Для этого он содержит ссылки на объект шины, набор подключённых хранилищ и модуль оптимизации внешнего распределения. Такая структура позволяет изолировать логику управления системой от конкретных реализаций хранения и передачи данных.

Основной задачей мастера является инициализация системы, которая включает подключение хранилищ к шине, настройку начального состояния распределённой среды и подготовку компонентов к выполнению запросов. Метод инициализации принимает конфигурационный объект, определяющий состав системы и параметры её работы.

После запуска системы мастер обеспечивает добавление и удаление узлов графа, направляя соответствующие операции в нужные хранилища через интерфейс шины. Благодаря этому достигается централизованное управление распределёнными изменениями при сохранении слабой связанности компонентов.

Метод оптимизации инициирует работу внешнего оптимизатора, который перераспределяет данные между хранилищами с целью улучшения локальности и уменьшения числа межшардовых переходов. Таким образом мастер выступает точкой запуска процедур улучшения распределения.

Дополнительно класс предоставляет метод поиска пути между вершинами, который реализует выполнение распределённого обхода графа с использованием шины и подключённых хранилищ. Результатом работы является строковое представление найденного пути, формируемое на основе последовательности посещённых вершин.

\begin{lstlisting}[language=C++, caption=Класс Master, label=impl-lst:master]
template<typename KeyType>
class Master {
private:
    Config cfg;
    IBus<KeyType>* bus;
    std::vector<IStorage<KeyType>*> storages;
    ExternalStorageOptimizer<KeyType>* optimizer = nullptr;
    IPathfinder pathfinder;

public:
    void set_bus(IBus<KeyType>* _bus);
    void set_optimizer(Optimizer<KeyType>* _optimizer);
    void add_storage(IStorage<KeyType>* storage);
    void remove_storage(int storage_id);
    IStorage<KeyType>* get_storage(int storage_id);
    const std::vector<IStorage<KeyType>*>& get_storages() const;

    void add_node(Node<KeyType> node);
    void inititalize(Config cfg);

    size_t storage_count() const;
    size_t total_nodes() const;
    size_t total_internal_edges() const;
    void run_optimization();

    std::string find_path(Node<KeyType> from, Node<KeyType> to);
};
\end{lstlisting}

\subsubsection{Классы распределения вершин}

\textbf{Реализация Fennel алгоритма}

Реализация потокового распределения в системе осуществляется за счёт специализированного расширения базового класса хранилища. Основная идея заключается в добавлении дополнительного состояния, необходимого для вычисления эвристики размещения вершин в рамках потокового алгоритма, а также в переопределении метода оценки изменения целевой функции при добавлении новой вершины.

В частности, расширенный класс содержит счётчик числа вершин в хранилище, который используется для оценки текущей нагрузки шарда и расчёта степени его заполненности. Данный параметр является ключевым при вычислении штрафных функций, связанных с балансировкой распределения.

Для работы алгоритма потокового распределения используется структура параметров \texttt{FennelParameters}, содержащая коэффициенты и настройки эвристики. Она позволяет задавать поведение алгоритма без изменения его кода, обеспечивая гибкость настройки модели распределения и возможность проведения экспериментов с различными конфигурациями.

Основной вычислительной частью реализации является функция оценки, основанная на модифицированной целевой функции Fennel. Данный метод учитывает как локальность размещения вершины относительно соседей, так и штраф за перегрузку шарда. В результате вычисляется значение, характеризующее изменение качества распределения при добавлении новой вершины в конкретное хранилище.

При добавлении вершины в хранилище происходит увеличение счётчика числа узлов, что позволяет поддерживать актуальное состояние нагрузки. Аналогично, при удалении вершины счётчик уменьшается, что обеспечивает корректное отражение текущего состояния системы в процессе динамического изменения графа.

Упрощённый вид изменений можно увидеть на листинге \ref{impl-lst:fennel}.

\begin{lstlisting}[language=C++, caption=Класс FennelStorage, label=impl-lst:fennel]
template <typename KeyType> 
class FennelStorage : public Storage<KeyType> {
protected:
FennelParameters params;
int number_of_nodes;

float get_p_parameter(int number_of_nodes, long number_of_edges);

float fennel_function(const Node<KeyType>& node, long total_edges);
public:
float get_streaming_euristics_change(const Node<KeyType>& node, long total_edges) override {
    return fennel_function(node, total_edges);
}

void get_add_announcement(Key key, int announcer_id, std::set<Edge<KeyType>> edges) override {
    ++number_of_nodes;
    Storage<KeyType>::get_add_announcement(key, announcer_id, edges);
};

void get_remove_announcement(Key key, int announcer_id) override {
    --number_of_nodes;
    Storage<KeyType>::get_remove_announcement(key, announcer_id);
}
};
\end{lstlisting}

\textbf{Реализация алгоритма уточнения Кернигана-Лина}

Реализация алгоритма Кернигана–Лина в подсистеме распределённого хранения данных осуществляется в рамках специализированного класса оптимизатора, основновные методы которого описаны в листинге \ref{impl-lst:KL}, предназначенного для перераспределения вершин графа между двумя хранилищами с целью уменьшения суммарного веса разреза и повышения локальности данных.

В отличие от компонентов подсистемы хранения, данный оптимизатор не является частью механизма взаимодействия между узлами системы и не реализует логику передачи данных. Его основная задача заключается в анализе текущего разбиения графа и вычислении эвристики улучшения распределения на основе структуры межшардовых связей. Для получения необходимой информации о границах разбиения используется интерфейс шины исключительно как источник данных о состоянии системы, без участия в процессе переноса или маршрутизации вершин.

Основой алгоритма является функция calculate\_gv, вычисляющая оценку изменения целевой функции для каждой вершины. Данная величина определяется как разность между суммарным весом внутренних и внешних рёбер вершины относительно текущего разбиения. Положительное значение соответствует предпочтительному размещению вершины в текущем хранилище, тогда как отрицательное значение указывает на потенциальное улучшение при её переносе.

На основе вычисленных значений формируется множество вершин с отрицательной оценкой, которые рассматриваются как кандидаты на перемещение между хранилищами. Эти вершины определяют направление итеративного улучшения разбиения графа.

Основной процесс оптимизации реализован в виде последовательного обмена вершинами между двумя хранилищами. На каждой итерации выбранные вершины удаляются из одного хранилища и добавляются в другое, после чего производится пересчёт эвристических значений с учётом изменённой структуры графа.

Итерационный процесс продолжается до достижения заданного ограничения по числу итераций, равного 5, так как это считается оптимальным числом при хорошем изначальном разбиении \cite{MetisOG}, либо до стабилизации состояния системы, при котором отсутствуют вершины, перемещение которых приводит к улучшению целевой функции.

\begin{lstlisting}[language=C++, caption=Класс KLExternalStorageOptimizer, label=impl-lst:KL]
template <typename KeyType>
class KLExternalStorageOptimizer {
private:
IBus<KeyType>* bus;

float calculate_gv(const Node<KeyType>& node, std::set<Edge<KeyType>> boundary_edges) const;
std::map<int, std::set<Node<KeyType>>> get_negative_gvs(std::map<int, std::map<Node<KeyType>, float>> full_map);
    
public:
std::map<int, std::map<Node<KeyType>, float>> calculate_gvs(int storage1, int storage2) const;
void optimize(int storage1, int storage2, int iterations_limit = 5);
};
\end{lstlisting}
